% Generated draft: requires mathematical and source review.
\documentclass[11pt]{article}
\usepackage[margin=25mm]{geometry}
\usepackage{fontspec}
\setmainfont{Latin Modern Roman}
\usepackage{amsmath,amssymb,mathtools,unicode-math}
\setmathfont{Latin Modern Math}
\usepackage{enumitem,needspace,xurl,hyperref}
\setlength{\parindent}{0pt}
\setlength{\parskip}{5pt}
\setlength{\emergencystretch}{3em}
\newcommand{\N}{\mathbb{N}}
\newcommand{\Z}{\mathbb{Z}}
\newcommand{\Q}{\mathbb{Q}}
\newcommand{\R}{\mathbb{R}}
\newcommand{\C}{\mathbb{C}}
\newcommand{\F}{\mathbb{F}}
\newcommand{\A}{\mathbb{A}}
\newcommand{\PP}{\mathbb{P}}
\newcommand{\E}{\mathbb{E}}
\newcommand{\eps}{\varepsilon}
\newcommand{\dd}{\,\mathrm{d}}
\begin{document}
\section*{UnsolvedMath ID: 30003840}
\textbf{Source problem code:} catalogue entry
\subsection{Definitions and mathematical statement}
Let \(F_{\mathrm{br}}\) and \(T_{\mathrm{br}}\) denote the braided Thompson groups intended in the cited Oberwolfach problem. The supplied record does not give presentations or constructions of these groups; their precise identification and effective input conventions therefore require confirmation from the source.

For a group \(G\), elements \(g,h\in G\) are conjugate if there exists \(x\in G\) such that \(x^{-1}gx=h\). In the usual finite-word formulation, fix a finite generating set \(S_G\), and write \([w]_G\) for the element represented by a word \(w\) over \(S_G\cup S_G^{-1}\). The conjugacy problem for \(G\) is decidable if there is an algorithm that, for every pair of such words \(u,v\), halts and correctly decides whether
\[
\exists x\in G\qquad x^{-1}[u]_Gx=[v]_G.
\]
Here the word-input convention explains the decision problem; the supplied record does not specify the generating sets or presentations.

Subject to the identification qualification above, the question is: for each \(G\in\{F_{\mathrm{br}},T_{\mathrm{br}}\}\), is its conjugacy problem decidable? The two groups are separate decision-problem targets; the question does not require a single common algorithm.
\subsection{Short English statement}
For each of the two braided Thompson groups named in the source, can an algorithm always decide whether two input words represent conjugate elements? The record leaves the concrete group definitions and word encodings unspecified. Its open-status assessment is supplied dated curation, not an independently verified conclusion.
\subsection{Sources}
\begin{itemize}
\item Local UnsolvedMath registry, supplied complete source record for permanent ID \(30003840\), titled \emph{Conjugacy in Braided Thompson Groups}; source problem code \texttt{catalogue entry}, in the Oberwolfach Reports Open Problems collection. The record gives proposed year 2018 and names no individual proposer.
\item Supplied report citation: \emph{Cohomological and Metric Properties of Groups of Homeomorphisms of R} (2019), Oberwolfach Reports, pp.~1579--1633, DOI \texttt{10.4171/owr/2018/26}. \url{https://doi.org/10.4171/owr/2018/26}. The record describes this collection as extracted from workshop open-problem discussions.
\item Supplied literature source labeled ``Oberwolfach Report 2019'': \url{https://ems.press/content/serial-article-files/46748}.
\item Supplied literature source labeled ``Bux--Flueck--Witzel'': \url{https://arxiv.org/abs/1210.2931}. No fuller bibliographic citation is supplied for this entry.
\item Literature assessment dated exactly \texttt{2026-08-21}: the supplied curation labels the problem open and reports that a 2018 result settles the conjugacy problem for \(V_{\mathrm{br}}\), explicitly leaves the corresponding questions for \(F_{\mathrm{br}}\) and \(T_{\mathrm{br}}\) open, and that no later general solution was found. This is supplied dated literature triage derived from the Oberwolfach curation, not independently verified here; its status and cited claims require independent verification.
\end{itemize}
\end{document}
